% beamer_template.tex
% 编译方式: xelatex beamer_template.tex

\documentclass[aspectratio=169, 12pt]{beamer}

% ============================================================
% 字体与中文支持
% ============================================================
\usepackage{fontspec}
\usepackage{xeCJK}
\setCJKmainfont[BoldFont=PingFang SC Semibold]{PingFang SC}
\setCJKsansfont[BoldFont=PingFang SC Semibold]{PingFang SC}
\setCJKmonofont{PingFang SC}
\setmainfont{Palatino}
\setsansfont{Helvetica Neue}
\setmonofont{Menlo}

% ============================================================
% Beamer 主题与颜色设置
% ============================================================
\usetheme{Madrid}
\usecolortheme{default}
\setbeamertemplate{navigation symbols}{} % 隐藏导航栏
\setbeamertemplate{footline}[frame number]

% ============================================================
% 常用宏包
% ============================================================
\usepackage{graphicx}
\graphicspath{{../key_slides/}}
\usepackage{xcolor}
\usepackage{hyperref}
\usepackage{booktabs}
\hypersetup{
  colorlinks=true,
  linkcolor=blue!70!black,
  urlcolor=blue!70!black,
}

% ============================================================
% 全局文档信息设置
% ============================================================
\title{Video Presentation Summary}
\subtitle{AI Generated Slides}
\author{Auto Generated by Video-to-Slides}
\date{\today}

\begin{document}

% ---------------------------------------
% 封面幻灯片
% ---------------------------------------
\begin{frame}[plain]
    \titlepage
\end{frame}

% ---------------------------------------
% 目录幻灯片
% ---------------------------------------
\begin{frame}{Contents}
    \tableofcontents
\end{frame}

% ============================================================
% 内容区域
% ============================================================

\section{Introduction}

\begin{frame}{First Slide Example}
    \begin{columns}
        % 左侧放文字要点
        \begin{column}{0.5\textwidth}
            \begin{itemize}
                \item Bullet point 1
                \item Bullet point 2
                \item Bullet point 3
            \end{itemize}
        \end{column}
        
        % 右侧放提取出来的图片
        \begin{column}{0.5\textwidth}
            \centering
            % \includegraphics[width=\textwidth]{example.jpg}
            \small Image Placeholder
        \end{column}
    \end{columns}
\end{frame}

\end{document}
